\subsection{Arithmetic normalization and a powered rule at every prime}
\label{sec:ftb-powered-key-all-primes}

The algebraic inputs here are projectivity and finite presentation of
the polarized abelian scheme; faithful-flat descent of its morphisms,
line bundles and identities; the origin-normalized theorem of the cube;
base change of its positive line-power section modules; and finite
etaleness of multiplication by an invertible integer.  These are the
same algebraic inputs as in the integral Koszul construction.  We keep
their applications explicit.

\begin{srlemma}[Countable descent preserving a local principal ideal base]
\label{lem:ftb-key-countable-descent}
Let \(R\) be a local principal ideal domain and \(x/R\) a key datum.
Every countable \(S\subset R\) is contained in a countable local
principal ideal subring \(R_0\subset R\), local and faithfully flat
in \(R\), over which \(x\) descends as a key datum.
Over a characteristic-zero field it descends to a countable subfield.
\end{srlemma}
\begin{proof}
The source constructs subrings as follows.
For countable \(S'\subset R\), put \(K=\operatorname{Frac}R\), let
\(K_0\) be the subfield generated by \(S'\), and set \(R_0=R\cap K_0\).
It is countable.  Divisibility descends: if \(b=ac\) with
\(a,b\in R_0\), then for \(a\ne0\), \(c=b/a\in R\cap K_0\);
the case \(a=0\) is immediate.  Units of \(R\) lying in \(R_0\)
therefore have their inverses in \(R_0\).
The subring is local and its map to \(R\) is local.
For an ideal \(I\subset R_0\), choose \(a\in I\) with \(aR\)
maximal among \(\{bR:b\in I\}\), using Noetherianity of \(R\).
Divisibility is total in \(R\).  For every \(b\in I\), maximality
then gives \(a\mid b\) in \(R\), hence in \(R_0\).
Thus \(I=aR_0\).  Torsion freeness over this principal ideal domain
makes \(R\) flat over \(R_0\); locality makes it faithfully flat.

Descend a projective presentation and its ample bundle using their
finite equations and gluing charts.  Descend the identity, inverse,
multiplication, and theta line bundle using finitely many further
coefficients.  Enlarge to a subring of the preceding kind.
The product used for multiplication is the descended product.
Faithful flatness tests its group identities and descends the
proper smooth polarized scheme, relative dimension, rigidified
symmetric ample line and principal polarization.
Thus the descended object is a polarized abelian scheme.

Over this local principal ideal ring, choose frames of the line at
the origin, the rank-one theta-section module and the rank-\(g\)
Lie module.  Their pullbacks differ from the specified data by an
actual line-bundle identification, a unit scalar and an invertible
matrix.  Enlarge \(S'\) to include these coefficients and the given
\(S\); the subring construction also includes the inverses of their
unit determinants.  Rescaling the frames gives the original datum
under the source's \(\mathrm{KeyPB}\), including its dual-number
Lie maps.  If \(R\) is a field, locality reflects every nonzero
element of \(R_0\) as a unit, so \(R_0\) is a subfield.
\end{proof}

\begin{srlemma}[Unit realization over a local DVR extension]
\label{lem:ftb-key-dvr-realization}
Fix admissible \(n,l,l',\alpha\), with
\(D=\prod_{j=0}^g(l-jn)(l'-jn)\).
For a key datum over a DVR \(R\) in which \(D\) is invertible,
there is an injective local map \(R\to R'\) to a DVR where
the exact \(\operatorname{KaVal}\) formula
\eqref{eq:ftb-kaform} is realized by a unit.
Its value commutes with field base change of that realization.
\end{srlemma}
\begin{proof}
Choose a basis of the level-\(n\) section module.
Over an algebraic closure of a rational-function extension of the
residue field choose \(g\) coefficient rows algebraically independent
over that field.  The preceding generic-section argument makes
their common zeros transverse.
Lift these finitely many coefficients after a local DVR extension
and completion.  Transcendental residue coefficients are obtained
by \(R[T_1,\ldots,T_b]_{(\pi)}\); further dominating DVR extensions
and completion permit the finitely many additional residue
coordinates.  This is the source's
\srcid{exists\_complete\_discreteValuationRing\_with\_residues}.
No quantitative extension degree is required.

Enlarge again to define the residue common zeros and the
\(D\)-torsion.  Multiplication by \(D\) is finite etale, and the
transverse zero scheme is etale at its zeros.  Hensel lifting over
the complete DVR supplies their sections with distinct reductions;
properness and the fixed zero count give all generic zeros.
The actual split-torsion and split-transverse-zero hypotheses
of Proposition~\ref{prop:ftb-integral-koszul-invariant} now hold.
It gives the unit and its termwise base-change identities.
The extension supplies a realization; it does not change the
section-independent field value.
\end{proof}

\begin{srproposition}[Rules away from one integer and field descent]
\label{prop:ftb-key-away-rule}
For each \(g\geq1\) and prime \(p\), there are positive integers
\(c,M\), \(p\nmid M\), and a section-independent rule on
characteristic-zero fields that is functorial for key base change,
satisfies
\[
                 a^{8c}\kappa_{\rm new}
                          =(\det B)^{4c}\kappa_{\rm old},
\]
has holomorphic normalized-period values for every
half-characteristic, and is a unit on each DVR where \(M\) is a unit.
Its unit specializes correctly if the residue field has
characteristic zero.
\end{srproposition}
\begin{proof}
Take \(n=4p,\ l=gn+1,\ l'=gn+1+p,\ M=D\);
Lemma~\ref{lem:ftb-key-numerical-data} gives \(p\nmid M\).
Proposition~\ref{prop:ftb-key-tuple-independent} defines a unique
nonzero value over algebraically closed fields and every generic
coefficient extension, agrees with every admissible tuple, and
gives field functoriality.  Its homogeneity is
\eqref{eq:ftb-koszul-homog}.
The preceding DVR lemma supplies a unit realization after a
local extension.  For a complex datum, descend to a countable
subfield, choose finitely many algebraically independent complex
coefficients, and obtain an actual transverse tuple.
Proposition~\ref{prop:ftb-key-holomorphic-period} supplies and
patches the local holomorphic value functions.

For an arbitrary characteristic-zero field \(K\), pull the datum
to \(\overline K\).  Every \(K\)-automorphism fixes its value by
functoriality and uniqueness.  An algebraic element with this
property lies in \(K\), since its minimal polynomial is separable.
This gives the unique descended nonzero value and its functoriality
under every field map.  To descend the homogeneity identity,
base change the two specified frames simultaneously and use
injectivity and the determinant base-change formula.

Let \(z\in\operatorname{Frac}(R)^\times\) be the descended value,
and \(u\in(R')^\times\) its unit realization.  In a common
algebraically closed fraction field their images agree.
Either \(z\in R\) or \(z^{-1}\in R\).
In the first case its image is \(u\), so the local map \(R\to R'\)
reflects it as a unit.  The second case applies the same argument
to \(z^{-1}\).  This descends the unit without requiring a finite
extension degree.  At a characteristic-zero residue map, pass to
a common algebraic closure of the residue fields of \(R'\) and
\(R\).  The actual determinant, cubical and Koszul formula
specializes there to the residue rule.  Locality makes the
residue maps factor through the original residue field; its
injective embedding in that closure proves the required
specialization equality.
\end{proof}

\begin{srlemma}[The precisely normalized elliptic test]
\label{lem:ftb-key-elliptic-discriminant}
For a smooth Weierstrass curve \(W/R\), the origin line bundle,
its canonical section and rigidification, and the tangent
\([\epsilon:-1:0]\), form a one-dimensional key datum.
For every complex normal form of its base change,
\[
                         \frac{M^4}{c_0^8}=s(\Delta_W).
 \tag{A1}\label{eq:ftb-key-elliptic}
\]
\end{srlemma}
\begin{proof}
The origin parameter on the dual-number tangent
\([\epsilon:-1:0]\) has value \(-\epsilon\).
The rigidified frame has value \(1\) at the origin, so the theta
function vanishes there and its derivative in that tangent
direction is \(-1\).
The three even genus-one half-characteristics give even theta
functions, by the substitution \(h\mapsto-h-2a\) in the series.
Their derivative is zero.  Hence the characteristic is the odd
one \((\tfrac12,\tfrac12)\).

Put \(b=2\pi i/M\).  The actual period lattice on that tangent line
is \(b(\mathbb Z+\tau\mathbb Z)\), and the derivative assertion says
\(\theta'[\tfrac12,\tfrac12](0)=-b/c_0\).
Jacobi's derivative formula in this normalization and the
Weierstrass period formula are
\[
 \theta'[\tfrac12,\tfrac12](0)^8=(2\pi)^8\Delta_{\rm mod}(\tau),
 \qquad
 \Delta\bigl(b(\mathbb Z+\tau\mathbb Z)\bigr)
                       =(2\pi/b)^{12}\Delta_{\rm mod}(\tau).
\]
Substitution gives the latter discriminant as
\((2\pi)^4b^{-4}c_0^{-8}=M^4/c_0^8\), since \(i^4=1\).
The actual Weierstrass uniformization identifies this lattice
discriminant with \(s(\Delta_W)\): its coordinate functions
satisfy the equation of \(W\), with the fixed origin parameter
\(t(0)=0,t'(0)=1\) and its prescribed leading coordinate limits.
These determine the same Weierstrass \(g_2,g_3\), equivalently
the standard periodic-solution uniqueness argument.
Thus both the tangent scale and the discriminant scale are fixed.
The analytic inputs are exactly the Jacobi derivative and
Weierstrass lattice identities displayed here.
\end{proof}

\begin{srproposition}[A normalized powered rule at every prime]
\label{prop:ftb-normalized-powered-key}
For each \(g\geq1\), a positive integer \(N(g)\) and a field-functorial
rule \(\rho_g\) exist with
\[
 \rho_g(x)=
       \left(\frac{(\det M)^4}{c_0^8}\right)^{N(g)}
 \tag{A2}\label{eq:ftb-powered-normal-value}
\]
in every complex half-characteristic normal form.
It obeys
\(a^{8N(g)}\rho_g(x_{\rm new})
=(\det B)^{4N(g)}\rho_g(x_{\rm old})\).
For every characteristic-zero DVR key datum its generic value
is a unit, with the correct specialization whenever the residue
field has characteristic zero.
\end{srproposition}
\begin{proof}
Choose an away-from-\(M\) rule in dimension \(2g\), exponent \(c>0\).
The finite-family argument gives the zero-characteristic constant
\(a_0\) of \eqref{eq:ftb-key-unscaled-constant}.
Evaluate this rule on \(x^2\).
Its doubled half-characteristic is even, and an integral symplectic
transformation takes it to zero.  The transformation preserves
\((\det M)^4/c_0^8\) by
\eqref{eq:ftb-normal-ratio-sp}; the square datum squares that ratio.
The resulting rule in dimension \(g\), exponent \(n=2c\), therefore
has normal value \(a_0((\det M)^4/c_0^8)^n\) for every
half-characteristic.  The field and DVR laws pass to the square.

For any integral Weierstrass curve \(W\), take its key datum over
\(\mathbb Z[\Delta_W^{-1}]\) and then its \(g\)-th power.
Lemma~\ref{lem:ftb-key-elliptic-discriminant} makes the value
\(a_0\Delta_W^{gn}\).  The field rule is defined over \(\mathbb Q\);
dividing its rational value by \(\Delta_W^{gn}\) proves
\(a_0=r\in\mathbb Q^\times\).
Whenever \(M\) and \(\Delta_W\) are DVR units, both the rule and
this discriminant power are units, so \(r\) is a unit there.
Use the two integral curves \(W_{j=0}\), with discriminant \(-27\),
and \(W_{j=1728}\), with discriminant \(-64\).
Both give the same \(r\), because they give the same complex \(a_0\).
In a local ring either \(3\) or \(1-3=-2\) is a unit.
One of these two discriminants is therefore a unit.
Thus \(r\) is a unit at every DVR where \(M\) is a unit.
Dividing by \(r\) gives exact normal value and preserves unit
and characteristic-zero specialization laws.

Do this separately for each prime \(p\), using \(M_p\) prime to
\(p\), and denote the normalized rule and exponent by
\((\rho_p,n_p)\).
For any two,
\[
                             \rho_p^{n_q}=\rho_q^{n_p}
 \tag{A3}\label{eq:ftb-powered-cross}
\]
on every characteristic-zero field.  Descend a datum to a
countable subfield and embed it in \(\mathbb C\); the equality
there is immediate from the exact normal values.
Functoriality and injectivity descend it.

Fix the rule at \(p=2\), setting \(N(g)=n_2\).
At a DVR where \(M_2\) is invertible it is already a unit.
Otherwise the residue field has a prime characteristic \(p\)
dividing \(M_2\).  The rule \(\rho_p\) is a unit there since
\(p\nmid M_p\); \eqref{eq:ftb-powered-cross} makes
\(\rho_2^{n_p}\) a unit.  Its valuation is
\(n_pv(\rho_2)\); positivity of \(n_p\) forces
\(v(\rho_2)=0\).
A characteristic-zero residue field makes every positive integer,
including \(M_2\), a unit; its specialization is therefore always
covered by the first case.  No excluded residue prime survives.
\end{proof}

\paragraph{Frozen source concordance and remaining producer.}
All paths have prefix
\srcid{src/Development/AbelianVarietyTheory/Proofs/ThetaLawProofs/}.
The local principal ideal subring is
\srcid{ExistsCountableLocalPidSubring.lean}
(\srcid{80f4e931782aab0999790894204dc2b851005194});
the actual scheme/group/line/frame descent is
\srcid{KeyDataDescendCountableSubrings.lean}
(\srcid{d06ec6cf161d65aff27cc78f9f3693252a1b58d2}).
DVR lifting and the away rule are
\srcid{ExistsKeyInvariantAfterDvrExtension.lean}
(\srcid{e9d5f9addb9796136f989270119a664641668639})
and \srcid{ExistsIntegralKeyInvariantRuleCharZero.lean}
(\srcid{86b60fd21a1e75dd48b4971ad05dc6e34e3f4228}).
Field descent and reflection of units are
\srcid{ExistsKeyPBPair.lean}
(\srcid{db42d4b4b8b9d8f9cd9a17237b338a72ed354076}).
The elliptic test is
\srcid{ExistsEllipticKeyDataValueEqDiscriminant.lean}
(\srcid{0e56730263497260189223973dbdaec4d38d1b09}).
Squaring, the rational tests and division by their common unit are
\srcid{ExistsNormalisedKeyInvariantRuleCharZero.lean}
(\srcid{e9c21ea730a06bb841cbfc137a7964f8f1c461c2}).
Cross powers and removal of all bad primes are
\srcid{ThetaKeyRuleExistsOfAlgebraicInvariants.lean}
(\srcid{24e1e9047fad6a46b1a8b2f4ecb81f4bbd835610}).

This is a powered rule with \(N(g)>0\), not yet the exponent-one
rule required by the key section.  Taking an arbitrary root would
lose field functoriality and specialization.  The source next fixes
that root by connected algebraic families with a known arithmetic
point.  The complex uniformization and polarized-torus
algebraization interfaces retained in the preceding subsection
also remain in force.
